% !PRINTED none
\weberpage{none}

\begin{center}
\gesperrt{Fünfundzwanzigster Abschnitt.}

\gesperrt{Transcendente Zahlen.}
\end{center}

\begin{center}
\rule{4cm}{0.4pt}
\end{center}

\begin{center}
\S.\,203.

\gesperrt{Abzählbare Mengen.}
\end{center}

In der Einleitung zu unserem Werke ist der allgemeine Begriff der Zahlen definirt worden, und mit diesem allgemeinen Zahlbegriffe haben wir zunächst, z.\,B. bei dem Beweise der Wurzelexistenz, operirt. Im weiteren Verlaufe unserer Betrachtungen haben wir uns dann nur mit \gesperrt{algebraischen Zahlen} beschäftigt, ohne uns darüber Rechenschaft zu geben, ob damit der Umfang des Zahlenbereiches erschöpft sei, oder ob es auch nichtalgebraische Zahlen giebt. Die Existenz von nichtalgebraischen Zahlen, die man auch transcendente Zahlen nennt, ist zuerst von \gesperrt{Liouville} bewiesen. Andere Beweise für diesen Satz rühren von \gesperrt{G. Cantor} her\footnote{\gesperrt{G. Cantor}, Crelle's Journal, Bd. 77 (1873). Ueber eine Eigenschaft des Inbegriffes aller reellen algebraischen Zahlen.}.

Wir knüpfen hier an den schon in der Einleitung aus einander gesetzten Begriff einer \gesperrt{Menge} oder \gesperrt{Mannigfaltigkeit} an, worunter ein System von Elementen irgend welcher Art zu verstehen ist, was so genau abgegrenzt ist, dass von jedem beliebigen Objecte völlig entschieden ist, ob es zu dem Systeme gehört oder nicht.

Wir unterscheiden endliche und unendliche Mengen, und führen als erstes und wichtigstes Beispiel einer unendlichen Menge die Gesammtheit der \gesperrt{natürlichen Zahlen} $1, 2, 3, \ldots$ an. Dann gilt folgende Definition:

1. \textbf{Definition.} \gesperrt{Eine Menge heisst abzählbar, wenn ihre Elemente mit der natürlichen Zahlenreihe}
%CONT>>
